%% ==========================================================================
%% Appendix D.3 — Murmuration Phase Transition Theory (Sections D.3.8–D.3.12)
%%
%% Drop-in LaTeX fragment for hjgauss.pdf, after Section D.3.7.
%% Follows the notation and style of hjgauss.pdf (NeurIPS 2026 submission).
%% Citations: CDC26 = \cite{Molu2026CDC} (CDC 2026 flocking hierarchy paper).
%% ==========================================================================

% --------------------------------------------------------------------------
\subsection{Vacuole Nucleation Topology}
\label{sec:d38_vacuole}
% --------------------------------------------------------------------------

When an attacking predator penetrates the interior of a flock's backward
reachable tube~(BRT), a topological hole---a \emph{vacuole}---nucleates in
the zero-sublevel set $\Omega^{-}(t) \triangleq \{x : v(t,x) \leq 0\}$.
We characterise this event via the Euler characteristic.

\begin{definition}[Euler Characteristic of the BRT]
\label{def:euler_char}
Let $\Omega^{-}(t)$ be the BRT at time~$t$, discretised on a uniform
$m\times m$ grid with vertex, edge, and face counts $V$, $E$, $F$.  The
\emph{Euler characteristic} is,
%
\begin{align}
\label{eq:euler_char}
\chi(t) \;\triangleq\; V - E + F \;=\; \beta_0(t) - \beta_1(t),
\end{align}
%
where $\beta_0$ is the number of connected components and $\beta_1$ is the
first Betti number (number of independent loops).
\end{definition}

\begin{theorem}[Vacuole Nucleation Topology]
\label{thm:vacuole_nucleation}
Let $v(t, x)$ be the viscosity solution of the backward reachable set (BRT)
equation~\eqref{eq:HJ-Visc}, and let $\Omega^{-}(t) = \{x \in \Omega : v(t,x) \leq 0\}$
denote the zero-sublevel set at time~$t$. Suppose the attacked-agent Hamiltonian
$H_{\mathrm{att}}(x, p)$ (Eq.~\eqref{eq:H_att}) satisfies $H_{\mathrm{att}}(x^*, p^*) = 0$
for some interior state $x^* \in \Omega$ with costate $p^* = Dv(t^*, x^*)$.

Then at time $t^*$ when this zero-crossing occurs, if the predator penetrates
the interior of the flock's reachable set, the topology of $\Omega^{-}(t^*)$
undergoes a transition characterized by,
%
\begin{align}
\label{eq:vacuole_threshold}
\chi(t^{*+}) \;=\; \chi(t^{*-}) - 1,
\end{align}
%
where $\chi$ is the Euler characteristic~\eqref{eq:euler_char}. This jump
indicates the nucleation of a vacuole (topological hole) in the flock's safe set,
marking the transition from simply connected ($\beta_1 = 0$) to multiply connected
($\beta_1 \geq 1$) topology.
\end{theorem}

\begin{proof}
\label{proof:vacuole_nucleation}
By Morse theory for sublevel sets of smooth functions, the topology of $\Omega^{-}(t)$
changes only at critical values of $v(\cdot, t)$.  Specifically, when the parameter
$t$ crosses a value $t^*$ at which a saddle point of $v(\cdot, t^*)$ lies on the
boundary of $\Omega^{-}(t^*)$, the sublevel set attaches a 1-handle.

The attachment of a 1-handle increases the first Betti number by one: $\beta_1(t^{*+}) = \beta_1(t^{*-}) + 1$.
Since the Euler characteristic is defined as $\chi = \beta_0 - \beta_1$ (Definition~\ref{def:euler_char}),
this increase in $\beta_1$ directly reduces $\chi$ by one:
%
\begin{align*}
\chi(t^{*+}) &= \beta_0(t^{*+}) - \beta_1(t^{*+}) \\
&= \beta_0(t^{*-}) - (\beta_1(t^{*-}) + 1) \\
&= \chi(t^{*-}) - 1.
\end{align*}
%
The condition $H_{\mathrm{att}}(x^*, p^*) = 0$ identifies the critical state-costate pair
where the attacked agent's optimal trajectory encounters the reachable set boundary.
At this instant, the saddle-point singularity of $v$ causes the topological change described above.
The geometric interpretation is that the predator has penetrated sufficiently deep into the flock
to create an enclosed region, manifested as a topological hole in the backward reachable set.
\end{proof}

% --------------------------------------------------------------------------
\subsection{Flock Splitting Dynamics}
\label{sec:d39_splitting}
% --------------------------------------------------------------------------

When a predator applies sufficient control to reduce the attacked agent's
optimal escape velocity, the BRT bifurcates into disjoint components.
This bifurcation corresponds to the moment when the attacking predator can
successfully isolate one or more agents from the main group.

\begin{proposition}[Bifurcation Condition for Flock Splitting]
\label{prop:bifurcation}
Let $\Omega^{-}(t)$ be the zero-sublevel set of the BRT at time~$t$.
The reachable set undergoes a bifurcation into $n_{\mathrm{comp}}(t) > 1$
connected components if and only if the attacked-agent Hamiltonian satisfies
the control-constrained minimax condition,
%
\begin{align}
\label{eq:bifurcation_condition}
H_{\mathrm{att}}(x^*, p^*) \;=\; \min_{u_p \in \mathcal{U}_p} \max_{u_e \in \mathcal{U}_e}
\left[
  p_1(u_p - v_e) + p_2 u_e + p_3(u_p^{z} - u_e^z) + p_4(\omega_p - \omega_e)
\right] \;<\; -\delta_c,
\end{align}
%
where $\delta_c > 0$ is a separation margin, $\mathcal{U}_p, \mathcal{U}_e$ are
admissible control sets, and $(u_p, \omega_p)$, $(u_e, \omega_e)$ are the
pursuer and evader acceleration and angular velocity controls respectively.
\end{proposition}

\begin{remark}[Temporal Bifurcation Pattern]
\label{rem:bifurcation_time}
The bifurcation occurs at a critical time $t_{\mathrm{split}}$ determined by
when the predator's control authority first permits the pursuer to violate
the costate constraint in the evader's optimal strategy. The sequence of
bifurcation times is monotone increasing: $t_{\mathrm{split}}^{(1)} < t_{\mathrm{split}}^{(2)} < \cdots$,
with the $m$-th bifurcation occurring when isolation of the $m$-th agent
becomes achievable under the game-theoretic equilibrium. The flock's resilience
to splitting depends on the coupling strength in the risk-reduction term
(the $p_4$ coefficient in the Hamiltonian), which represents the predation
risk shared across connected components.
\end{remark}

% --------------------------------------------------------------------------
\subsection{Cordon Formation and Capture Prevention}
\label{sec:d310_cordon}
% --------------------------------------------------------------------------

A flock executing a defensive \emph{cordon formation} creates an annular
perimeter around threatened agents, preventing isolated predator incursions into
the interior. This configuration is geometrically characterized by a multiply
connected reachable set with one independent loop ($\beta_1 = 1$).

\begin{proposition}[Annular BRT and Barrier Certificate]
\label{prop:cordon}
Let $\mathcal{C}(r_{\mathrm{in}}, r_{\mathrm{out}}) \triangleq \{x \in \Omega : r_{\mathrm{in}}
\leq \|x_{1:2}\| \leq r_{\mathrm{out}}\}$ denote the annular barrier region (``cordon'')
with inner radius $r_{\mathrm{in}}$ and outer radius $r_{\mathrm{out}}$.
If the BRT restricted to $\mathcal{C}$ satisfies,
%
\begin{align}
\label{eq:cordon_euler}
v(t, x) \;>\; 0 \quad \text{for all } x \in \{y \in \mathcal{C} : \|y_{1:2}\| < r_{\mathrm{in}}\},
\end{align}
%
then the inner region is a \emph{barrier-certificate-protected domain}, and any trajectory of
the attacked agent passing through $\mathcal{C}(r_{\mathrm{in}}, r_{\mathrm{out}})$ remains captured
(i.e., reaches the terminal set) within finite time. The topology of $\Omega^{-}(t) \cap \mathcal{C}$
is annular, with Betti number $\beta_1 = 1$ and Euler characteristic,
%
\begin{align}
\label{eq:cordon_euler_char}
\chi\bigl(\Omega^{-}(t) \cap \mathcal{C}\bigr) \;=\; 0.
\end{align}
\end{proposition}

\begin{remark}[Barrier Certificate Interpretation]
\label{rem:barrier_cert}
The annular topology (Eq.~\eqref{eq:cordon_euler_char}) physically represents
a mutual-defense configuration where the safe set (valued region, $v \leq 0$)
wraps around the predator or protected interior, creating a topological
obstruction to escape. From the reachability perspective, any agent trapped
inside the cordon cannot escape; from outside, no agent can cross the barrier
without incurring catastrophic cost. This geometric property is independent
of the detailed dynamics and depends only on the level-set structure of $v$.
\end{remark}

The numerical validation of the cordon configuration is visualised through
reachability level sets (files \texttt{reachability\_t\{:\!04d\}.jpg} in
the simulation output directory), which show the barrier topology explicitly.

% --------------------------------------------------------------------------
\subsection{Flash Expansion Events}
\label{sec:d311_flash}
% --------------------------------------------------------------------------

The \emph{flash expansion}~(IJRR23 Action~4) is characterised by a rapid,
isotropic outward motion of the flock boundary in response to a predator
approach.  In the BRT framework, this corresponds to a linear growth of the
zero-level-set radius.

\begin{lemma}[Linear Flash Expansion Rate]
\label{lem:flash_rate}
Under the 4D aerial dynamics~\eqref{eq:rel_dynamics_4d}, the maximum radial
extent of the BRT zero-level-set in the horizontal plane grows at the rate,
%
\begin{align}
\label{eq:flash_rate}
\max_i \lVert x_i(t+1)_{1:2} \rVert
- \max_i \lVert x_i(t)_{1:2} \rVert
\;=\; \Theta\!\left(T_2^{-1}\right),
\end{align}
%
where $T_2$ is the tipping time of the second hierarchy level~(CDC26
Lemma~2).  The altitude component $x_3$ does not contribute to the radial
growth rate because the capture set is a cylinder~\eqref{eq:terminal_cost_4d}.
\end{lemma}

\begin{proof}
\label{proof:flash_rate}
By the comparison principle for viscosity solutions, $v(t,x) \leq 0$ implies
$v(t+\tau, x + \tau f(x,u)) \leq 0$ for any admissible trajectory.  The
maximum horizontal speed of the free-agent dynamics is $v_e$ (linear speed
bound), giving a radius growth of at most $v_e \Delta t$ per step.  The lower
bound $\Omega(T_2^{-1})$ follows from CDC26 Lemma~1: below the tipping time
$T_2$, the pursuer cannot reduce the evader's speed below a constant fraction
of $v_e$, ensuring the growth rate is bounded away from zero.
\end{proof}

The 4D flash-expansion rate converges to the 3D rate as altitude variance
$\sigma^2_{x_3} \to 0$, consistent with the altitude-decoupling property of
the cylinder terminal cost (see Test~\texttt{test\_altitude\_decoupling} in
\S\ref{sec:d312_markers}).

% --------------------------------------------------------------------------
\subsection{Phase Transition Markers in Numerical Solutions}
\label{sec:d312_markers}
% --------------------------------------------------------------------------

We detect phase transitions computationally by monitoring the topological
invariants $\chi(t)$, $\beta_1(t)$, and $n_{\mathrm{comp}}(t)$ of the
zero-sublevel set across backward time steps.  The detection algorithm is:

\begin{enumerate}
  \item At each step~$t$, compute the 2D slice $v(t, x_1, x_2;\bar{x}_3, 0)$
        on a grid of resolution $m \times m$.
  \item Label connected components of $\{v \leq 0\}$ via
        \textsc{scipy.ndimage.label}.
  \item Estimate $\beta_1$ by counting encircled complementary components.
  \item Compute $\chi(t) = n_{\mathrm{comp}}(t) - \beta_1(t)$.
\end{enumerate}

The threshold-crossing rules that trigger an event label are:
%
\begin{align}
\label{eq:vacuole_marker}
\Delta\chi_t \;\triangleq\; \chi_{t+1} - \chi_t \;<\; -0.5
&\quad\Longrightarrow\quad \text{vacuole nucleation,} \\
\label{eq:cordon_marker}
\beta_1(t) \;\geq\; 1 \;\text{ and }\; \beta_1(t-1) = 0
&\quad\Longrightarrow\quad \text{cordon formation,} \\
\label{eq:frag_marker}
n_{\mathrm{comp}}(t) \;>\; n_{\mathrm{comp}}(t-1)
&\quad\Longrightarrow\quad \text{flock fragmentation.}
\end{align}

The phase-transition classification is validated against CDC26 Table~I, which
partitions the state space into the five canonical murmuration phases: (i)
cohesion, (ii) evasion, (iii) cordon, (iv) expansion, and (v) fragmentation.
Each detected event is appended to \texttt{phase\_transitions.txt} with its
time index, event type, $\chi$, $\beta_1$, and $n_{\mathrm{comp}}$ values.
Figure~\ref{fig:topology_summary} shows the topology metrics $\chi(t)$,
$\beta_1(t)$, and $n_{\mathrm{comp}}(t)$ over time, generated by the
\texttt{OutputHandler.save\_topology\_summary} routine (file
\texttt{/tmp/murmurations/topology\_summary.jpg}).
